\Article{19}
If we resolve the magnetic force at the point $O$ into three mutually perpendicular components $X,Y,Z$, with the third directed toward the center of the Earth, then $X$ and $Y$ are tangent to the spherical surface that passes through $O$ and is concentric with the Earth. Specifically, $X$ is directed northward in the plane through $O$ and the Earth's axis, and $Y$ westward parallel to the Earth's equator. We then have
\[
X=-\frac{\mathrm{d}V}{r\mathrm{d}u},\quad Y=-\frac{\mathrm{d}V}{r\sin u\,\mathrm{d}\lambda},\quad Z=-\frac{\mathrm{d}V}{\mathrm{d}r}
\]
and consequently
\[
\begin{aligned}
X&=-\frac{R^3}{r^3}\left(\frac{\mathrm{d}P'}{\mathrm{d}u}+\frac Rr\cdot\frac{\mathrm{d}P''}{\mathrm{d}u}+\frac{RR}{rr}\cdot\frac{\mathrm{d}P'''}{\mathrm{d}u}+\text{etc.}\right)\\
Y&=-\frac{R^3}{r^3\sin u}\left(\frac{\mathrm{d}P'}{\mathrm{d}\lambda}+\frac Rr\cdot\frac{\mathrm{d}P''}{\mathrm{d}\lambda}+\frac{RR}{rr}\cdot\frac{\mathrm{d}P'''}{\mathrm{d}\lambda}+\text{etc.}\right)\\
Z&=\frac{R^3}{r^3}\left(2P'+\frac{3RP''}{r}+\frac{4RRP'''}{rr}+\text{etc.}\right)
\end{aligned}
\]

On the Earth's surface, $X,Y$ are the same horizontal components denoted above by these letters, and $Z$ is the vertical component, positive when directed downward. The expressions for these forces on the Earth's surface are therefore
\SourceLeaf{147}
\[
\begin{aligned}
X&=-\left(\frac{\mathrm{d}P'}{\mathrm{d}u}+\frac{\mathrm{d}P''}{\mathrm{d}u}+\frac{\mathrm{d}P'''}{\mathrm{d}u}+\text{etc.}\right)\\
Y&=-\frac1{\sin u}\left(\frac{\mathrm{d}P'}{\mathrm{d}\lambda}+\frac{\mathrm{d}P''}{\mathrm{d}\lambda}+\frac{\mathrm{d}P'''}{\mathrm{d}\lambda}+\text{etc.}\right)\\
Z&=2P'+3P''+4P'''+\text{etc.}
\end{aligned}
\]

\Article{20}
Let us now combine these propositions with the known theorem that every function of $\lambda$ and $u$ which has a definite finite value for all values of $\lambda$ from $0$ to $360^0$, and of $u$ from $0$ to $180^0$, can be expanded in a series of the form
\[
P^0+P'+P''+P'''+\text{etc.}
\]
whose general term $P^{\mathrm{n}}$ satisfies the partial differential equation above, that such an expansion is possible in just one definite way, and that this series always converges. We obtain the following remarkable propositions.

I. Knowing the value of $V$ at all points on the Earth's surface is sufficient to derive from it the general expression for $V$ throughout the infinite space outside the Earth's surface, and thus also to determine the forces $X,Y,Z$ on the Earth's surface and throughout the infinite space outside it. Clearly, all that is needed is to expand $\frac VR$ in a series according to the theorem just mentioned.

In what follows, therefore, the symbol $V$ is always to be understood in the sense restricted to the Earth's surface, unless the contrary is expressly stated, or as the function of $\lambda$ and $u$ obtained from the general expression by setting $r=R$, so that
\[
V=R(P'+P''+P'''+\text{etc.})
\]

II. Knowing the value of $X$ at all points on the Earth's surface is sufficient to obtain everything stated in I. Indeed, according to Art. 15, we have $\int_0^u X\mathrm{d}u=\frac{V^0-V}{R}$, where $V^0$ denotes the value of $V$ at the north pole. The expansion of $\int_0^u X\mathrm{d}u$ in a series of the stated form must necessarily be identical with
\SourceLeaf{148}
\[
V^0-P'-P''-P'''-\text{etc.}
\]

III. In the same way, with reference to Art. 16, it is clear that knowing the value of $Y$ over the entire Earth, together with knowing $X$ at all points on a line running from one pole of the Earth to the other, suffices to establish the \textit{complete} theory of terrestrial magnetism.

IV. Finally, it is clear that the complete theory can also be derived simply from knowledge of the value of $Z$ over the entire surface of the Earth. Indeed, if $Z$ is expanded in a series
\[
Z=Q^0+Q'+Q''+Q'''+\text{etc.}
\]
whose general term satisfies the partial differential equation mentioned several times above, then we must have $Q^0=0$ and $P'=\frac12Q'$, $P''=\frac13Q''$, $P'''=\frac14Q'''$ and so on.

\Article{21}
Because the individual forces $X,Y,Z$ depend in such a simple way on a single function $V$ and are so simply related to each other, they are far better suited to serve as the basis of the theory than the usual expression of magnetic force by the three elements of total intensity, inclination, and declination. Or rather, the latter form. natural as it may seem in itself when all that matters is to grasp the facts, cannot directly serve to establish the theory, at least in the first instance, until it is converted into the other form. In this respect, it would therefore be very desirable to have a general graphical representation of the horizontal intensity. In part this is because it is closer than the total intensity to what the theory can use, and in part because, in the great majority of cases, it was the horizontal intensity that was actually observed in the first place, whereas the total intensity was only calculated from it using the inclination. Keeping the elements of horizontal magnetism separate is all the more advisable because the means now available allow them to be determined with markedly greater precision. At the very least, the total intensity calculated from the observed horizontal \SourceLeaf{149}intensity should never be published while withholding the observed value unless the inclination used in the calculation is also given, so that anyone who wants to use the data for the theory can recover the original figures without distortion.

It would certainly be interesting to base the entire theory of terrestrial magnetism on observations of the horizontal needle alone and thereby anticipate the vertical component or the inclination. For now, however, it is much too soon for that. The shortcomings of the data currently available do not allow us to dispense with using the vertical component as well. Indeed, the theory already gains confirmation if it can be shown that all the elements are compatible under a single principle.


